\documentclass{article}
\usepackage[T1]{fontenc}
\usepackage{lmodern}
\usepackage{graphicx}
\usepackage{amsmath,amssymb}
\usepackage[a4paper,margin=2cm]{geometry}
\usepackage[absolute,overlay]{textpos}
\setlength{\TPHorizModule}{1mm}
\setlength{\TPVertModule}{1mm}
\usepackage[indonesian]{babel}
\usepackage{hyperref}
\setlength{\emergencystretch}{2em}
% unit: Halaman 29

\begin{document}

% === Halaman 29 (llm) ===

\textbf{Contoh 9.8.} Tinjau kembali plainteks dari Contoh 9.6:

\begin{center}
$10100010001110101001$
\end{center}

Bagi plainteks menjadi blok-blok yang berukuran 4 bit:

\begin{center}
$1010 \quad 0010 \quad 0011 \quad 1010 \quad 1001$
\end{center}

atau dalam notasi HEX adalah $\text{A23A9}$.

Misalkan kunci ($K$) yang digunakan adalah (panjangnya juga 4 bit)

\begin{center}
$1011$
\end{center}

atau dalam notasi HEX adalah $\text{B}$. Sedangkan $IV$ yang digunakan seluruhnya bit 0 (Jadi, $C_0 = 0000$)

Fungsi enkripsi $E$ yang digunakan sama seperti sebelumnya: XOR-kan blok plainteks $P_i$ dengan $K$, kemudian geser secara \textit{wrapping} bit-bit dari $P_i \oplus K$ satu posisi ke kiri.

\[ E_K(P) = (P \oplus K) \ll 1 \]

\end{document}
